\title[Prime-to-$p$ sections separate rational points]{Rational points of curves in characteristic $p$ are separated by the prime-to-$p$ fundamental group, on a question of Grothendieck}
\author{}
\date{23 November 2026}

\begin{document}

\begin{abstract}
Grothendieck observed, in a manuscript now in the Montpellier archive, that over a field of characteristic $p$ the abelianized prime-to-$p$ fundamental group of a curve does not separate its rational points, because prime-to-$p$ Kummer theory does not see torsion of $p$-power order, and he asked whether the whole prime-to-$p$ fundamental group does. For hyperbolic curves over finite fields, Sa\"idi and Tamagawa proved that the decomposition groups in this group separate the closed points outside a finite exceptional set. We prove that there is no exception. Let $X$ be a hyperbolic curve over a field $K$ finitely generated over $\F_p$. Then two distinct rational points of $X$ never define conjugate sections of the geometrically prime-to-$p$ fundamental group of $X$ over the Galois group of $K$. When $K$ is finite it follows that the exceptional set of Sa\"idi and Tamagawa is empty: every closed point is determined by its decomposition group. The proof rests on the comparison theorems of SGA~1, on the Riemann hypothesis for curves over finite fields, on the finiteness theorem of Sa\"idi and Tamagawa, and on the configuration space used in Mochizuki's theory of cuspidalization. The theorem and its proof were produced by GPT-6 Astra (ChatGPT Pro mode) and checked by Claude Opus 5.5 at effort max. Claude Fable 5.2 at effort max edited the text, followed every step again and added the corollary on closed points, which has had no second check. No human has checked any of this.
\end{abstract}

\maketitle

\section{The question and the result}

\subsection{Grothendieck's question}
Let $K$ be a field, $\overline K$ a separable closure, $G_K=\mathrm{Gal}(\overline K/K)$, and $X$ a smooth geometrically connected curve over $K$. The fundamental group of $X$ is an extension of $G_K$ by the fundamental group of $X_{\overline K}$, and a rational point $P\in X(K)$ splits it: $P$ gives a section $s_P$ of $\pi_1(X)\to G_K$, well defined up to conjugation by $\pi_1(X_{\overline K})$. The first question of the arithmetic of fundamental groups is whether $P\mapsto[s_P]$ is injective. Over a finitely generated field of characteristic $0$ it is, for curves of positive genus already in the abelianization: the difference of two sections is there the Kummer class of $P-Q$ in the Jacobian, and by the Mordell--Weil theorem this class vanishes only if $P=Q$.

In characteristic $p$ the fundamental group of $X_{\overline K}$ is replaced by its maximal quotient of order prime to $p$, the quotient that specializes well and is topologically finitely generated. In a manuscript of the Montpellier archive \cite[cote 149, ms.~pp.~17--18]{Archive}, Grothendieck notes that the abelian argument then fails, already for an elliptic curve with a point removed: the prime-to-$p$ Kummer map of the Jacobian kills its $p$-primary torsion. He asks whether the full prime-to-$p$ group may still separate the points. He reduces the question to a statement about coverings: if $P$ and $Q$ have the same class, then for every connected prime-to-$p$ Galois cover and every point $P'$ above $P$ there is a point $Q'$ above $Q$ such that $Q'-P'$ has $p$-power order in the generalized Jacobian of the cover. He writes that one might be able to deduce $P=Q$ from this.

\subsection{The theorem}
Let $p$ be a prime, $K$ a field finitely generated over $\F_p$, and $X$ a smooth geometrically connected curve over $K$ which is \emph{hyperbolic}: $2g-2+n>0$, where $g$ is the genus of the smooth compactification of $X_{\overline K}$ and $n$ the number of its points at infinity. Put
\[
 \Delta_X=\pi_1(X_{\overline K})^{(p')},\qquad
 \Pi_X=\pi_1(X)\big/\ker\bigl(\pi_1(X_{\overline K})\to\Delta_X\bigr),
\]
where $(p')$ denotes the maximal quotient of order prime to $p$. The sequence $1\to\Delta_X\to\Pi_X\to G_K\to1$ is exact: only the geometric part has been made prime to $p$, and the arithmetic quotient is the whole of $G_K$. We call $\Pi_X$ the \emph{geometrically prime-to-$p$ fundamental group} of $X$. A point $P\in X(K)$ defines a section $s_P$ of $\Pi_X\to G_K$, up to conjugation by $\Delta_X$; its image is a decomposition group of $P$. We write $\Sec(\Pi_X\to G_K)$ for the set of continuous sections.

\begin{theorem}\label{thm:main}
The map
\[
 X(K)\longrightarrow \Sec(\Pi_X\to G_K)/\Delta_X,\qquad P\longmapsto[s_P],
\]
is injective.
\end{theorem}

The theorem says that two distinct rational points can be told apart by a finite \'etale cover. If $P\ne Q$, there is a finite \'etale cover $Y\to X$, whose pullback to $X_{\overline K}$ is dominated by a Galois cover of degree prime to $p$, with a $K$-rational point above $P$ and none above $Q$. Indeed, the covers of this kind with a rational point above $P$ correspond to the open subgroups of $\Pi_X$ containing $s_P(G_K)$; if each of them had a rational point above $Q$, the limit argument of Step~4 below would give a decomposition group of $Q$ inside $s_P(G_K)$, hence equal to it. In Grothendieck's form: if $P,Q\in X(K)$ satisfy his condition for all prime-to-$p$ covers, then $P=Q$; for proper curves over a finite field the two forms are equivalent by Kummer theory, since the kernel of the prime-to-$p$ Kummer map of a finite abelian group is its $p$-primary part.

The hypothesis that $X$ be hyperbolic cannot be dropped. For an elliptic curve $E$ the group $\Delta_E$ is abelian, and $s_P$ and $s_Q$ are conjugate exactly when $P-Q$ has $p$-power order in $E(K)$. Removing a point from $E$ does not change the abelianization of the geometric group, which is the content of Grothendieck's remark; by Theorem~\ref{thm:main} it makes the map injective. For $\mathbb A^1$ and $\mathbb P^1$ the group $\Delta_X$ is trivial and the map is constant. For $\mathbb G_m$ the map is injective, by the Kummer argument of Section~\ref{sec:open}.

Over a finite field every closed point becomes rational after a finite extension of the base field, and the theorem then gives a statement about all closed points. Let $k$ be a finite field of characteristic $p$, $X$ a hyperbolic curve over $k$, and $\widetilde X$ the normalization of the smooth compactification of $X$ in the pro-cover of $X$ defined by $\Pi_X$. For a closed point $\tilde x$ of $\widetilde X$ let $D_{\tilde x}\subset\Pi_X$ be its decomposition group, the stabilizer of $\tilde x$. Sa\"idi and Tamagawa call $\tilde x$ \emph{exceptional} if there is a closed point $\tilde x'\ne\tilde x$ such that $D_{\tilde x}\cap D_{\tilde x'}$ is open in both groups; the images in $X$ of the exceptional points form their exceptional set \cite[Definitions 1.3, 1.5, 1.10]{ST09}.

\begin{maincor}\label{cor:closed}
Let $X$ be a hyperbolic curve over a finite field of characteristic $p$. No closed point of $\widetilde X$ is exceptional. In particular, the decomposition groups in $\Pi_X$ of two distinct closed points of $X$ are not conjugate.
\end{maincor}

\subsection{Why it matters}
For the full fundamental group of a hyperbolic curve over a finite field, the map from closed points to conjugacy classes of decomposition groups is injective, and the anabelian theorems of Tamagawa and Mochizuki recover the curve from the group; Sa\"idi and Tamagawa recall this in \cite[Introduction]{ST09}. For the geometrically prime-to-$p$ group, the natural one in families and the one of Grothendieck's question, they proved that the map is injective outside a finite exceptional set \cite[Proposition 1.8(vi), Corollary 1.11]{ST09}, wrote that they could only prove it ``almost injective'', and recovered from the group the points of the complement of that set \cite[Introduction]{ST09}. In their later refinement for geometrically pro-$\Sigma$ groups the exceptional sets could not even be shown to be finite \cite[Introduction, pp.~4--5]{ST18}. Theorem~\ref{thm:main} answers Grothendieck's question for rational points of hyperbolic curves over every finitely generated field of characteristic $p$, and Corollary~\ref{cor:closed} shows that for the set of all primes different from $p$ the exceptional set is empty.

\subsection{Credit}
Grothendieck asked the question and reduced it to the statement about coverings \cite[ms.~pp.~17--18]{Archive}. Step~4 uses the finiteness of the exceptional set, which is a theorem of Sa\"idi and Tamagawa \cite[Proposition 1.8(vi)]{ST09}. Step~5 uses Lemma 1.7 of the same paper, and Corollary~\ref{cor:closed} uses parts (ii) and (iii) of its Proposition 1.8. The objects of Steps~1, 2 and~5 are those of Mochizuki's theory of cuspidalization \cite{Moch}: the configuration space of two points on the curve, the decomposition group of the diagonal, and the quotients of the fundamental group of a punctured curve that are pro-$r$ or central along the puncture, whose extension class is the Chern class of the point. Sa\"idi and Tamagawa carried that theory over to prime-to-$p$ groups \cite[\S 2]{ST09}; we cite it from their account, and we do not use its reconstruction theorems. The comparison and specialization theorems are those of SGA~1 \cite{SGA1}. The absolute values of the Frobenius eigenvalues on $H^1$ of a curve are given by Weil's Riemann hypothesis for curves, cited here from \cite[Th\'eor\`eme 1.6]{WeilI}.

What is new here is the following. Lemma~\ref{lem:free} and Proposition~\ref{prop:exact} prove directly that the configuration fibration stays exact when it is made pro-$r$ along the diagonal. A conjugator between the sections of $P$ and $Q$ is then used to carry the curve punctured at $P$ to the curve punctured at $Q$, over the identity of $\Pi_C$ (Step~2). Steps~3 and~4 compare the numbers of points of corresponding covers and deduce that the points outside the exceptional set are fixed. Step~5 is a Kummer argument with the function whose divisor is $p^e(P-Q)$.

\subsection{Context}
In characteristic $0$ the injectivity of $P\mapsto[s_P]$ follows from the Mordell--Weil theorem, as recalled above, and the open part of the section conjecture is surjectivity. In characteristic $p$, Grothendieck's remark shows why the prime-to-$p$ form is harder: the abelian information is blind to the $p$-power torsion of Jacobians, so that an argument must use the nonabelian structure of $\Delta_X$. The theorem of Sa\"idi and Tamagawa does so, through cuspidalization, and leaves the finite exceptional set. We did not find Theorem~\ref{thm:main} or Corollary~\ref{cor:closed} in the papers cited here. We could not consult Stix's monograph \cite{Stix}, which has a chapter on the injectivity of $P\mapsto[s_P]$, and so we cannot say what it contains on prime-to-$p$ quotients in characteristic $p$.

\subsection{Scope}
Proved: Theorem~\ref{thm:main}, for rational points of proper and of affine hyperbolic curves over every field finitely generated over $\F_p$; and Corollary~\ref{cor:closed}, for all closed points when the base field is finite. Not proved: the analogue for geometrically pro-$\Sigma$ quotients, $\Sigma$ a smaller set of primes; any statement on surjectivity, that is, on which sections come from points; any statement in higher dimension. For an infinite field $K$ we have not written out the statement about decomposition groups of closed points of higher degree.

\section{Proof overview}

The case to which everything is reduced is that of a proper curve $C$ of genus at least two over a finite field $k$. Suppose that $P\ne Q$ in $C(k)$ have conjugate sections in $\Pi_C$. Then $P-Q$ is torsion of $p$-power order in the Jacobian, and this is exactly what prime-to-$p$ Kummer theory cannot see. The difficulty is to turn the equality of two sections of $\Pi_C$, which says nothing about functions, into a statement about a function on $C$. The idea is to carry the equality from $C$ to the punctured curves $C-P$ and $C-Q$, whose fundamental groups do know the functions with zeros and poles at $P$ and $Q$. This is done in three moves.

First, the configuration space $C\times C$ minus the diagonal fibres over $C$, and its fibre over $P$ is $C-P$. Fix a prime $r\ne p$. When the fundamental groups are made pro-$r$ along the diagonal and along the puncture, the fibre sequence stays exact (Step~1), and the decomposition group of the diagonal contains the full inertia group $\Z_r(1)$ of the puncture (Step~2). A conjugator between $s_P$ and $s_Q$ lifts to the decomposition group of the diagonal. Conjugation by the lift carries the fibre over $P$ to the fibre over $Q$, and gives an isomorphism $\beta$ between these groups of $C-P$ and of $C-Q$ which induces the identity on $\Pi_C$ and respects the punctures. Genus at least two is used here: the fibre group must be free of rank at least two for Lemma~\ref{lem:free}. The restriction to prime-to-$p$ groups is used to compare with a lift to characteristic $0$ and to have tame ramification along the diagonal.

Second, $\beta$ preserves the $r$-adic cohomology of the compactifications of corresponding covers, with its Frobenius action. By the trace formula, corresponding covers have the same number of points over every finite extension of $k$ (Step~3). A limit argument then shows that $\beta$ sends the decomposition group of a rational point to the decomposition group of a rational point. Since $\beta$ is the identity on $\Pi_C$, the two points have the same decomposition group in $\Pi_C$, and outside the finite exceptional set of Sa\"idi and Tamagawa they are the same point (Step~4).

Third, let $f$ be the function with divisor $p^e(P-Q)$. On the quotients that are central along the punctures, $\beta$ defines a class with residues $1$ at $P$ and $-1$ at $Q$. The Kummer class of $f$ differs from $p^e$ times this class by the Kummer class of a constant $a\in k^\times$, because $H^1$ of the curve with coefficients $\Z_r(1)$ has no Frobenius invariants: this is where the Riemann hypothesis for curves enters. At a point $R$ outside the exceptional set the class defined by $\beta$ vanishes, so $f(R)/a$ has order prime to $r$. But $f$ takes every value, among them $a\zeta$ with $\zeta\ne1$ a root of unity of order a power of $r$ (Step~5).

The general case follows in Section~\ref{sec:open}. A proper curve over a finitely generated field is specialized to a closed point of a model where $P$ and $Q$ stay distinct. An affine curve of genus at least two is mapped to its compactification. In genus $0$ and $1$ the abelianization suffices, after one prime-to-$p$ cover in genus $1$; these cases use that $K$ is finitely generated, through the theorem of Lang and N\'eron and the freeness of divisor groups. Corollary~\ref{cor:closed} is deduced from Theorem~\ref{thm:main} over the finite extensions of the base field.

